\chapter{Generative Representations}
\label{ch:generative-representations}

\emph{Opening dilemma.} The course-support team has 200 labeled URGENT messages and 18{,}000 ROUTINE ones. Their triage model collapses to ``always predict routine'' before training even finishes. They consider three responses. Hire annotators to label more urgent cases --- expensive and slow. Oversample the 200 they already have --- the model memorizes the rare class instead of generalizing. Or generate synthetic urgent messages from a model trained on the few they have --- cheap, fast, and potentially worthless if the generated messages just average together what the real urgent messages had in common. The chapter's running question is the last one: can we generate examples that genuinely add information about the underrepresented region of the data space, and how would we know the synthetic ones are not just elaborate memorization?

The previous chapter treated representations as tools for \emph{understanding}---spaces that make downstream tasks easier. Representations can also \emph{generate}: produce new data points that resemble the training data but were never observed. A course-support system that generates realistic synthetic student messages can augment scarce labeled data. A pedestrian detector that generates novel nighttime scenes can fill gaps in its training distribution. The same latent spaces that enable generation also enable interpolation, anomaly detection, and controllable editing.

This chapter introduces three families of generative models---variational autoencoders, generative adversarial networks, and diffusion models. Each learns a latent representation from which new data can be sampled \citep{kingma2014vae,goodfellow2014gan,ho2020ddpm}. The probabilistic foundations from Chapter~\ref{ch:probabilistic-models} motivate VAEs and latent-variable sampling; GANs introduce adversarial training. One practical question runs through all three: how do you know if generated data is any good?

Generation is not an end in itself. Every generative model encodes assumptions about which variations matter and which should be preserved. The design question is not ``can the model generate?'' but ``does what it generates serve the downstream task, and how would we know?''

\textbf{Representation for generation.} When evaluating a generative model, ask what the latent space \emph{means}. Can you move through the latent space and predict how the output will change? If not, the model may still generalize, but the latent space becomes harder to trust for controlled generation or interpretation.

\begin{practicalnote}
\textbf{Depth guide.} In a first course, the core expectation is to compare generative families, explain their objectives at a design level, and evaluate their downstream risks. Training stable GANs, deriving diffusion objectives in full, or building a production generator belongs in extended projects.
\end{practicalnote}

\section{From Representation to Generation}

Chapter~\ref{ch:representation-foundation-models} introduced autoencoders: networks that compress an input into a bottleneck representation and then reconstruct it. The bottleneck forces the network to preserve the most important structure in a compact code. A standard autoencoder's latent space, however, has no particular structure---it is not organized for sampling new points. Pick a random point in the latent space of a trained autoencoder, decode it, and the result is usually meaningless noise.

Generative models impose structure on the latent space so it \emph{can} be sampled. The three families differ in how they achieve this:

\begin{itemize}[leftmargin=*]
\item \emph{Variational autoencoders (VAEs)} regularize the latent space into a smooth, continuous distribution. New points can then be sampled from it and decoded into realistic outputs.
\item \emph{Generative adversarial networks (GANs)} learn a mapping from a simple random distribution to realistic outputs by training a generator to fool a discriminator.
\item \emph{Diffusion models} learn to reverse a gradual noising process, starting from pure noise and iteratively denoising it into a data sample.
\end{itemize}

Each family trades off sample quality, diversity, training stability, and computational cost differently.

\section{Variational Autoencoders}
\label{sec:vae}

\subsection{The Idea: Structured Latent Spaces}

A variational autoencoder (VAE) modifies the standard autoencoder in one key way: the encoder maps each input not to a single latent point but to a \emph{distribution} over latent points. The encoder typically outputs the mean $\mu$ and per-coordinate variance $\sigma^2$ of a diagonal Gaussian in latent space.

This design matters because the VAE does not optimize reconstruction alone. Training also applies a regularization term that pulls each input's latent distribution toward a shared prior, usually a standard Gaussian. This KL pressure prevents the encoder from scattering training examples into isolated, disconnected pockets of latent space. Instead, it produces a latent space that is smooth and organized: nearby latent points decode to similar outputs, and gradual movement through the latent space produces gradual changes in the reconstruction.

Once the latent space has this structure, sampling becomes meaningful rather than arbitrary. During training, a latent point is sampled from the encoder's distribution and passed to the decoder. After training, new data can be generated by sampling from the prior. Interpolating between two latent codes also produces a smooth transition between the corresponding outputs.

\subsection{The VAE Objective}

A VAE is trained on a single objective that balances two goals:

\textbf{Reconstruction.} The decoded output should match the original input. As with a standard autoencoder, a reconstruction loss measures this---mean squared error for continuous data, cross-entropy for discrete data.

\textbf{Regularization.} The encoder's output distribution should stay close to a simple prior, typically a standard Gaussian $\mathcal{N}(0, I)$. The KL divergence measures the gap between the two:
\[
D_{\text{KL}}\bigl(q(z \mid x) \;\|\; p(z)\bigr)
\]
where $q(z \mid x) = \mathcal{N}(\mu(x), \operatorname{diag}(\sigma^2(x)))$ is the encoder's output and $p(z) = \mathcal{N}(0, I)$ is the prior.

The total loss is:
\[
\mathcal{L}_{\text{VAE}} = \underbrace{\mathbb{E}_{q(z \mid x)}\bigl[-\log p(x \mid z)\bigr]}_{\text{reconstruction}} + \underbrace{D_{\text{KL}}\bigl(q(z \mid x) \;\|\; p(z)\bigr)}_{\text{regularization}}
\]

This is the negative \emph{evidence lower bound} (ELBO) used in variational autoencoders \citep{kingma2014vae,rezende2014sgvb}. Minimizing it improves reconstruction while keeping the latent space organized.

\begin{prereqbridge}
\textbf{Why we cannot just maximize $\log p(x)$.} The marginal $p(x) = \int p(x \mid z) p(z)\, dz$ integrates over every latent code that could have produced $x$. For all but trivial $p(x\mid z)$ this integral has no closed form and is too high-dimensional to evaluate by quadrature, so direct gradient ascent on $\log p(x)$ is not an option. The ELBO replaces that intractable objective with a tractable lower bound; the slack between the two is exactly the gap between the approximating posterior $q(z\mid x)$ and the true posterior.
\end{prereqbridge}

\textbf{Why the ELBO lower-bounds $\log p(x)$ (skip on a first pass).} The argument is one application of Jensen's inequality (Mathematical Bridge II). The marginal log-likelihood of the data is
\[
\log p(x) \;=\; \log \int p(x \mid z)\, p(z)\, dz.
\]
For any approximating distribution $q(z \mid x)$ with full support on the relevant $z$, multiply and divide by $q(z \mid x)$ inside the integral:
\[
\log p(x) \;=\; \log \mathbb{E}_{q(z \mid x)}\!\left[\frac{p(x \mid z)\, p(z)}{q(z \mid x)}\right]
\;\ge\; \mathbb{E}_{q(z \mid x)}\!\left[\log\frac{p(x \mid z)\, p(z)}{q(z \mid x)}\right],
\]
where the inequality is Jensen's: $\log$ is concave, so $\log \mathbb{E}[Z] \ge \mathbb{E}[\log Z]$. Splitting the right-hand side into reconstruction and regularization terms gives
\[
\mathrm{ELBO}(x; \theta, \phi) \;=\; \underbrace{\mathbb{E}_{q(z \mid x)}[\log p(x \mid z)]}_{\text{reconstruction (negate this term)}} \;-\; \underbrace{D_{\mathrm{KL}}\bigl(q(z \mid x)\,\|\,p(z)\bigr)}_{\text{regularization}}.
\]
The slack is exactly $\log p(x) - \mathrm{ELBO} = D_{\mathrm{KL}}\bigl(q(z \mid x) \,\|\, p(z \mid x)\bigr) \ge 0$, with equality only when the approximating posterior matches the true one. Maximizing the ELBO with respect to both the encoder $\phi$ and the decoder $\theta$ therefore pushes $\log p(x)$ up while pushing $q$ toward the true posterior. The VAE training loss is $-\mathrm{ELBO}$, which is the expression at the top of the section.

\textbf{Closed-form KL for diagonal Gaussians.} For $q(z \mid x) = \mathcal{N}(\mu, \mathrm{diag}(\sigma^2))$ and $p(z) = \mathcal{N}(0, I)$ in $d$ dimensions,
\[
D_{\mathrm{KL}}\bigl(q \,\|\, p\bigr) \;=\; \tfrac{1}{2}\sum_{j=1}^{d}\bigl(\mu_j^2 + \sigma_j^2 - \log \sigma_j^2 - 1\bigr).
\]
The formula is what the implementation actually computes; expanding the Gaussian densities and integrating term by term recovers it. The reconstruction term, $-\mathbb{E}_{q(z|x)}[\log p(x \mid z)]$, has no closed form in general and is estimated by a one-sample Monte Carlo at each gradient step---which is what the reparameterization trick below makes practical.

For one input, the encoder might output \(\mu=(0.4,-0.2)\) and \(\log\sigma^2=(-1.0,0.2)\). The model samples noise \(\epsilon\sim\mathcal{N}(0,I)\), forms \(z=\mu+\sigma\odot\epsilon\), decodes \(z\), and measures the reconstruction loss. For a diagonal Gaussian, the KL term reduces to
\[
\frac{1}{2}\sum_j \left(\mu_j^2+\sigma_j^2-\log\sigma_j^2-1\right).
\]
Large means or extreme variances inflate this penalty. The encoder is therefore rewarded for choosing a latent distribution that reconstructs the input while staying close enough to the shared prior for sampling to remain meaningful.

A weighting factor $\beta$ controls the balance between reconstruction and regularization (the ``$\beta$-VAE'' variant). At $\beta = 1$, the objective is the standard ELBO; even then, the two terms need not be equal in scale or contribution. Larger $\beta$ raises the KL pressure, often producing a more structured, disentangled latent space at the cost of blurrier reconstructions. Smaller $\beta$ does the opposite: sharper reconstructions, less-organized latent space. In practice, tuning $\beta$ is often necessary to get useful results for the downstream task.

\subsection{The Reparameterization Trick}

Training a VAE requires backpropagating gradients through the sampling step, but sampling is not differentiable. The \emph{reparameterization trick} fixes this by rewriting the sample as a deterministic transformation of a noise variable:
\[
z = \mu(x) + \sigma(x) \odot \epsilon, \quad \epsilon \sim \mathcal{N}(0, I)
\]

The randomness now lives entirely in $\epsilon$, which is independent of the model parameters. Gradients flow through $\mu$ and $\sigma$ back to the encoder. This trick makes VAE training practical with standard backpropagation.

The whole VAE training step fits in twenty lines of numpy-style pseudocode. The next snippet wraps a forward pass, the reparameterized sample, the decoder, the closed-form Gaussian KL, and the per-batch loss the autodiff system would differentiate.

\begin{codeexample}
import numpy as np

def vae_step(x, encoder, decoder, rng, beta=1.0):
    """One forward pass of a VAE: returns loss and its components.
    encoder(x) -> (mu, log_sigma2), each of shape (batch, latent_dim).
    decoder(z) -> x_recon of the same shape as x (Gaussian-decoder mean).
    """
    # 1. Encode the input to the variational posterior.
    mu, log_sigma2 = encoder(x)
    sigma          = np.exp(0.5 * log_sigma2)
    # 2. Reparameterized sample: z = mu + sigma * eps, with eps ~ N(0, I).
    eps  = rng.standard_normal(size=mu.shape)
    z    = mu + sigma * eps
    # 3. Decode and compute Gaussian negative log-likelihood
    #    (here written as squared error, which is -log N(x | x_recon, I) + const).
    x_recon = decoder(z)
    recon_term = 0.5 * np.sum((x - x_recon) ** 2, axis=-1)        # per example
    # 4. Closed-form diagonal-Gaussian KL.
    kl_term    = 0.5 * np.sum(mu ** 2 + sigma ** 2 - log_sigma2 - 1.0, axis=-1)
    # 5. ELBO loss = -ELBO = reconstruction NLL + beta * KL, averaged over batch.
    loss = (recon_term + beta * kl_term).mean()
    return loss, recon_term.mean(), kl_term.mean()
\end{codeexample}

The five-step structure is the algorithm. In a real PyTorch or JAX implementation, every named variable is a tensor with autograd tracking enabled, and the framework's autodiff produces gradients of the loss with respect to all encoder and decoder parameters in one backward pass. The reparameterized $z = \mu + \sigma \odot \epsilon$ is the line that makes that backward pass possible: $\epsilon$ has no gradient because it is sampled from a fixed distribution, so the gradient of any function of $z$ flows directly into $\mu$ and $\sigma$ through their algebraic dependencies.

\subsection{What VAEs Are Good For}

\textbf{Latent interpolation.} Given two inputs $x_1$ and $x_2$, encode them and interpolate between their posterior means, for example $z_t = (1-t)\mu(x_1) + t\mu(x_2)$, then decode. A well-trained VAE produces smooth transitions: morphing between two faces, shifting between writing styles, or blending pedestrian poses.

\textbf{Anomaly detection.} Inputs unlike the training data tend to produce high reconstruction error, low model likelihood or ELBO, or an unusual posterior relative to the learned aggregate posterior. This makes VAEs useful for detecting out-of-distribution inputs, though the score must be validated in the deployment domain---a capability that connects directly to the epistemic uncertainty discussion in Chapter~\ref{ch:probabilistic-models}.

\textbf{Data augmentation.} Sampling the latent space generates new training examples. Unlike rule-based augmentation (Chapter~\ref{ch:dataset-design}), VAE-based augmentation produces novel examples rather than transforming existing ones. The same validation discipline still applies: augmented data must be evaluated against a real held-out test set.

\begin{example}[VAE for Student Message Augmentation]
The course-support team has 200 labeled URGENT messages but needs more to train a reliable classifier. They train a VAE on all 12,000 student messages, ignoring labels. The VAE learns a latent space where similar messages cluster together. The team encodes the 200 URGENT messages, samples new latent points near their encodings (with small added noise), and decodes them. The generated messages resemble real URGENT messages without being copies. After manual inspection removes implausible generations, the team adds 150 synthetic URGENT messages to the training set. Recall on URGENT messages improves from 0.71 to 0.79 on the held-out gold test set.
\end{example}

\begin{warningbox}
VAE-generated samples tend to be blurrier than real data, especially for images. The Gaussian decoder assumption and the reconstruction loss, which averages over possible outputs, are responsible. For tasks where sharp, realistic detail matters---generating training images for a pedestrian detector, for example---GANs or diffusion models often produce better samples. For tasks where diversity and coverage matter more than sharpness, such as generating text variations, VAEs are usually sufficient.
\end{warningbox}

\begin{failurebox}
\textbf{Posterior collapse.} A subtler VAE failure than blurriness is \emph{posterior collapse}: the encoder learns to ignore the input entirely and outputs the prior $q(z \mid x) \approx p(z)$ for every $x$, regardless of what $x$ is. The KL term in the ELBO is then zero (the encoder matches the prior exactly), and the decoder learns to do everything from the unconditional prior --- in effect, the VAE becomes an unconditional generator with the encoder doing nothing. This is especially common with very expressive decoders (e.g., autoregressive RNNs or transformers as the decoder $p(x \mid z)$), because such decoders can model $p(x)$ well without needing any latent information. Diagnose it by watching the KL term during training: if it drops sharply to near zero and stays there, the posterior has collapsed. The standard mitigations are \emph{KL annealing} (ramp $\beta$ from $0$ up to $1$ during training so the reconstruction term anchors $z$ first), \emph{free bits} (skip the KL gradient when it falls below a per-dimension floor), and constraining the decoder so it cannot model $p(x)$ unconditionally well.
\end{failurebox}

\section{Generative Adversarial Networks}
\label{sec:gan}

\emph{Back to the running dilemmas.} VAE samples for the course-support team are smooth, easy to interpolate, and visibly blurry --- a synthetic ``urgent'' message reads like an average of urgency rather than any particular case. The pedestrian-detector team has the same complaint about VAE-generated nighttime frames: the silhouettes are placed correctly, but the textures lack the sharp detail a downstream detector needs to learn from. GANs sit at the opposite end of the trade: sharp, photorealistic samples produced by a generator that has been trained against a discriminator, with all the training-stability costs that ``two networks chasing each other'' implies.

\subsection{The Adversarial Idea}

Generative adversarial networks (GANs) take a different approach. Instead of learning a latent distribution explicitly, a GAN learns to generate data through a two-player game:

\textbf{The generator} $G$ takes a random noise vector $z \sim \mathcal{N}(0, I)$ and produces a synthetic data point $G(z)$.

\textbf{The discriminator} $D$ takes a data point---real or generated---and outputs the probability that it is real.

The two networks train simultaneously with opposing objectives. The discriminator tries to distinguish real from generated data. The generator tries to fool the discriminator. Formally:
\[
\min_G \max_D \; \mathbb{E}_{x \sim p_{\text{data}}}[\log D(x)] + \mathbb{E}_{z \sim p(z)}[\log(1 - D(G(z)))]
\]
This is the original minimax GAN objective from \citet{goodfellow2014gan}; practical GAN variants modify the loss or regularization to improve stability.

At equilibrium, in theory, the generator produces data indistinguishable from real data, and the discriminator outputs 0.5 for everything.

\textbf{The optimal discriminator and the JS-divergence connection (Extension).} Fix the generator $G$ and ask which $D$ maximizes the inner objective. Writing $p_g$ for the distribution of generated samples, the discriminator's expected payoff is
\[
V(D) \;=\; \int \bigl[p_\text{data}(x)\,\log D(x) + p_g(x)\,\log(1 - D(x))\bigr]\,dx.
\]
Pointwise this is $a \log d + b \log(1 - d)$ with $a = p_\text{data}(x)$, $b = p_g(x)$, $d = D(x)$; setting the derivative to zero gives the optimum $d^\star = a/(a+b)$, so
\[
D^\star(x) \;=\; \frac{p_\text{data}(x)}{p_\text{data}(x) + p_g(x)}.
\]
Substituting this back into the generator's objective and rewriting the integrand using the average distribution $m(x) = \tfrac{1}{2}(p_\text{data}(x) + p_g(x))$ gives, up to an additive constant,
\[
\max_D V(G, D) \;=\; -\log 4 \;+\; 2 \cdot D_\mathrm{JS}\bigl(p_\text{data} \,\|\, p_g\bigr),
\]
where $D_\mathrm{JS}(p \,\|\, q) = \tfrac{1}{2} D_\mathrm{KL}(p \,\|\, m) + \tfrac{1}{2} D_\mathrm{KL}(q \,\|\, m)$ is the Jensen--Shannon divergence. So with an optimal discriminator the generator is, in principle, minimizing JS divergence to the data distribution. Two things follow. First, the global optimum is $p_g = p_\text{data}$, at which point $D^\star \equiv \tfrac{1}{2}$ everywhere and the loss equals $-\log 4$. Second, when the supports of $p_g$ and $p_\text{data}$ are disjoint (a common case early in training when the generator is bad), $D_\mathrm{JS}$ saturates at $\log 2$ and its gradient with respect to the generator's parameters \emph{vanishes}: the discriminator can perfectly separate real from fake, but the JS landscape is locally flat for the generator. This is the formal reason GAN training stalls when the discriminator is too good, and the motivation for the Wasserstein GAN, which replaces JS with the Earth-Mover (Wasserstein-1) distance whose gradient remains informative even under disjoint support.

\subsection{Why GANs Produce Sharp Samples}

VAEs optimize a reconstruction loss that averages over uncertainty, producing blurry outputs. GANs instead use a discriminator that penalizes any detectable difference from real data. This adversarial pressure pushes the generator toward sharp, realistic outputs. The visual difference is striking: GAN-generated faces can be nearly photorealistic, while VAE-generated faces are typically smoother and less detailed.

\subsection{Training Instability}

Adversarial training is inherently unstable. The generator and discriminator must stay in balance. If the discriminator becomes too strong, every generated sample is detected and the generator gets no useful gradient signal. If the generator becomes too strong too quickly, the discriminator stops providing meaningful feedback. Common failure modes include:

\textbf{Mode collapse.} The generator produces only a few types of outputs that fool the discriminator, ignoring the diversity of real data. A face-generating GAN might produce only young female faces; a pedestrian-generating GAN might produce only frontal-view standing poses. The discriminator sees realistic samples but cannot detect the lack of diversity.

\textbf{Oscillation.} The two networks cycle: the generator finds an exploit, the discriminator catches on, the generator switches to a new exploit. Training loss oscillates without converging.

\textbf{Vanishing gradients.} An overconfident discriminator (outputs near 0 or 1) provides almost no gradient signal to the generator, stalling training.

All three failures share a single mathematical cause: when $p_\text{data}$ and $p_g$ have disjoint or nearly-disjoint support, $\mathrm{JSD}(p_\text{data} \,\|\, p_g) \to \log 2$ regardless of how far apart they are. The discriminator can perfectly distinguish them ($D^\star$ saturates at $0$ or $1$), the gradient $\nabla_g \log(1 - D(G(z)))$ vanishes wherever $D$ is saturated, and the generator receives no useful signal about which direction to move. Mode collapse, oscillation, and vanishing gradients are three behavioral symptoms of the same geometric problem: JS-divergence is constant once the supports separate, and so is the gradient. The Wasserstein-GAN fix replaces JS-divergence with the Earth Mover's distance, which scales with how far apart the supports are even when they are disjoint, restoring a meaningful gradient signal.

\begin{warningbox}
Mode collapse is insidious because standard training metrics may not detect it. The discriminator loss can converge, the generator loss can decrease, and individual samples can look realistic---yet the \emph{distribution} of generated samples may cover only a fraction of the real distribution. Always inspect the diversity of generated samples, not just their quality.
\end{warningbox}

Several techniques improve GAN training stability:
\begin{itemize}[leftmargin=*]
\item \emph{Wasserstein GAN (WGAN):} Replaces the original objective with the Wasserstein distance, which provides smoother gradients even when the discriminator is strong.
\item \emph{Spectral normalization:} Constrains the discriminator's Lipschitz constant, preventing it from becoming too sharp.
\item \emph{Progressive growing:} Starts at low resolution and gradually increases it, giving the generator time to learn coarse structure before fine detail.
\item \emph{Two-timescale learning rates:} Uses a higher learning rate for the discriminator to keep it slightly ahead of the generator and provide useful gradients.
\end{itemize}
None of these fully solves the stability problem. GAN training demands more monitoring and hyperparameter tuning than standard supervised training.

\subsection{Conditional Generation}

A \emph{conditional GAN} (cGAN) adds a conditioning input---a class label, a text description, or an image---so the generator produces outputs that match the condition. The generator becomes $G(z, c)$, the discriminator becomes $D(x, c)$, and $c$ is the condition.

\begin{example}[Conditional GAN for Pedestrian Scenes]
A pedestrian detection team wants to augment its training data with nighttime scenes---an epistemic gap identified in the uncertainty audit. They train a conditional GAN where the condition is ``time of day'' (day, dusk, night). Given the condition ``night,'' the generator produces dark scenes with streetlight illumination and pedestrian silhouettes. After adding the synthetic nighttime scenes to the training set, the detector's nighttime recall improves from 0.62 to 0.74 on real nighttime test frames. Separate evaluation on the real test set confirms the gap between synthetic and real.
\end{example}

\section{Diffusion Models}
\label{sec:diffusion}

\emph{Why the pedestrian-detector team will use this section.} VAEs produced smooth but soft synthetic frames; GANs produced sharp ones but with mode collapse and training pain. Diffusion sits at a third corner of the trade: training is regression-stable, sample quality typically exceeds GANs on natural images, and conditioning (``pedestrian crossing a wet road at dusk'') is straightforward. The cost is inference latency --- many forward passes per sample --- which matters whether the synthetic frames are for offline data augmentation (latency is tolerable) or for real-time use (it isn't). The chapter's running dilemma about whether synthetic data adds information becomes more answerable here, because diffusion's mode coverage tends to be better than GANs.

\subsection{The Denoising Idea}

Diffusion models are a leading family in image and media generation, building on denoising diffusion and score-based generative modeling \citep{ho2020ddpm,song2021score}. Their core idea is surprisingly simple: learn to reverse a gradual noising process.

\textbf{Forward process (noising).} Start with a real data point $x_0$. Over $T$ steps, add small amounts of Gaussian noise to produce a sequence $x_1, x_2, \ldots, x_T$. For $T$ large enough, $x_T$ is indistinguishable from pure Gaussian noise.

\textbf{Reverse process (denoising).} Train a neural network $\epsilon_\theta(x_t, t)$ to predict the noise added at each step. To generate new data, start from pure noise $x_T \sim \mathcal{N}(0, I)$ and iteratively denoise: at each step, the network predicts and removes the noise, gradually revealing a realistic data point.

The training objective is simple. Minimize the mean squared error between the predicted noise and the actual noise added at each step:
\[
\mathcal{L}_{\text{diffusion}} = \mathbb{E}_{t, x_0, \epsilon}\bigl[\|\epsilon - \epsilon_\theta(x_t, t)\|^2\bigr]
\]

\textbf{The forward process has a closed form.} Each forward step adds Gaussian noise according to a chosen variance schedule $\beta_t \in (0, 1)$:
\[
q(x_t \mid x_{t-1}) \;=\; \mathcal{N}\bigl(x_t;\; \sqrt{1 - \beta_t}\, x_{t-1},\; \beta_t I\bigr).
\]
Writing $\alpha_t = 1 - \beta_t$ and $\bar\alpha_t = \prod_{s=1}^{t} \alpha_s$, the marginal at step $t$ given the original $x_0$ collapses, by closure of Gaussians under convolution, to a single Gaussian:
\[
q(x_t \mid x_0) \;=\; \mathcal{N}\bigl(x_t;\; \sqrt{\bar\alpha_t}\, x_0,\; (1 - \bar\alpha_t)\, I\bigr),
\quad\text{equivalently}\quad
x_t \;=\; \sqrt{\bar\alpha_t}\, x_0 + \sqrt{1 - \bar\alpha_t}\, \epsilon,\; \epsilon \sim \mathcal{N}(0, I).
\]
This is what makes diffusion training tractable. Rather than unrolling $t$ noising steps one at a time, we sample $t$ uniformly, sample one $\epsilon$, form $x_t$ in one shot, and ask the network to predict the $\epsilon$ that produced it. The MSE objective above is a reweighted variational bound on $\log p(x_0)$; choosing the noise-prediction parameterization and dropping the per-$t$ weights produces a clean training signal that empirically works better than the strict ELBO version.

\begin{advancednote}
\textbf{Denoising is score matching: the Tweedie identity.} Tweedie's formula states that for $x \sim \mathcal{N}(\mu, \sigma^2 I)$, the posterior mean of $\mu$ given $x$ is
\[
\mathbb{E}[\mu \mid x] \;=\; x + \sigma^2 \nabla_x \log p(x).
\]
Apply this to the diffusion forward marginal $x_t \mid x_0 \sim \mathcal{N}(\sqrt{\bar\alpha_t}\, x_0,\; (1 - \bar\alpha_t) I)$. Reading off $\mu = \sqrt{\bar\alpha_t}\, x_0$ and $\sigma^2 = 1 - \bar\alpha_t$, and using the reparameterization $x_t = \sqrt{\bar\alpha_t}\, x_0 + \sqrt{1 - \bar\alpha_t}\, \epsilon$, a few lines of algebra give the optimal noise predictor
\[
\epsilon^\star(x_t, t) \;=\; -\sqrt{1 - \bar\alpha_t}\, \nabla_{x_t} \log p_t(x_t),
\]
where $p_t$ is the marginal of $x_t$ under the forward process.

This is the identity that makes ``noise prediction'' (used in DDPM training) and ``score matching'' (used in score-based generative models, Song and Ermon 2019) two views of the same object. Training a network $\epsilon_\theta(x_t, t)$ to predict the noise $\epsilon$ added in the forward process is equivalent --- up to the fixed scaling $-\sqrt{1 - \bar\alpha_t}$ --- to learning the score $\nabla_{x_t} \log p_t(x_t)$ of the noisy marginal. The reverse-time SDE that turns noise back into data only needs the score, so either parameterization plugs into the same sampler.

The implication is conceptual, not just notational: denoising-diffusion and score-based generative modeling are not parallel ideas, but the same algorithm written in different parameterizations.
\end{advancednote}

\begin{example}[The Forward Chain by Hand]
Take a scalar $x_0 = 2.0$, a three-step schedule $\beta_1 = 0.1,\, \beta_2 = 0.2,\, \beta_3 = 0.3$, and three independent noise draws $\epsilon_1 = 0.5,\, \epsilon_2 = -0.3,\, \epsilon_3 = 1.2$ from $\mathcal{N}(0, 1)$. The iterative form $x_t = \sqrt{1 - \beta_t}\, x_{t-1} + \sqrt{\beta_t}\, \epsilon_t$ walks the chain one step at a time:
\begin{align*}
x_1 &= \sqrt{0.9}\,(2.0) + \sqrt{0.1}\,(0.5)  \approx 1.897 + 0.158 = 2.056,\\
x_2 &= \sqrt{0.8}\,(2.056) + \sqrt{0.2}\,(-0.3) \approx 1.839 - 0.134 = 1.704,\\
x_3 &= \sqrt{0.7}\,(1.704) + \sqrt{0.3}\,(1.2)   \approx 1.426 + 0.657 = 2.083.
\end{align*}
The closed form skips the intermediates. Compute $\bar\alpha_3 = 0.9 \cdot 0.8 \cdot 0.7 = 0.504$, so
\[
x_3 \;=\; \sqrt{0.504}\,(2.0) + \sqrt{0.496}\,\epsilon \;\approx\; 1.420 + 0.704\,\epsilon,\quad \epsilon \sim \mathcal{N}(0, 1).
\]
A single $\epsilon$ here is the convolution of the three step noises, weighted by the geometry of the chain; the iterative draw above happens to correspond to $\epsilon \approx 0.94$. Both forms give the same marginal $q(x_3 \mid x_0) = \mathcal{N}(1.420,\, 0.496)$, and the closed form is what training actually samples from.

The diffusion forward process is a deterministic shrinkage of the data signal by $\sqrt{\bar\alpha_t}$ combined with progressively dominant Gaussian noise of variance $1 - \bar\alpha_t$. By step $T$, $\bar\alpha_T \to 0$, the signal is gone, and $x_T$ is approximately $\mathcal{N}(0, 1)$---which is exactly where the reverse process starts.
\end{example}

\textbf{The reverse step has a closed form too (Extension).} A first-pass reader can take the DDPM sampling code below as the operational rule and skip the derivation; the construction here is for readers who want to see where the update comes from. A short Bayes-rule calculation gives the posterior $q(x_{t-1} \mid x_t, x_0)$ as Gaussian with mean $\tilde\mu_t(x_t, x_0)$ depending on $x_t$ and $x_0$. Substituting the predictor $\epsilon_\theta(x_t, t)$ for $\epsilon$ via $x_0 \approx (x_t - \sqrt{1-\bar\alpha_t}\,\epsilon_\theta)/\sqrt{\bar\alpha_t}$ gives the DDPM sampling update:
\[
x_{t-1} \;=\; \frac{1}{\sqrt{\alpha_t}}\!\left(x_t - \frac{1 - \alpha_t}{\sqrt{1 - \bar\alpha_t}}\,\epsilon_\theta(x_t, t)\right) \;+\; \sigma_t\, z,\quad z \sim \mathcal{N}(0, I),
\]
with $\sigma_t^2$ chosen from the same $\beta$-schedule (typically $\sigma_t^2 = \beta_t$ or the smaller $\tilde\beta_t = \frac{1-\bar\alpha_{t-1}}{1-\bar\alpha_t}\beta_t$). At $t = 1$ the noise term is dropped to get a deterministic final denoising. The sampling loop runs this update for $t = T, T-1, \ldots, 1$ starting from $x_T \sim \mathcal{N}(0, I)$; DDIM's deterministic variant sets $\sigma_t = 0$ and skips intermediate $t$'s, which is why it produces high-quality samples in $10$--$50$ steps rather than the original $T = 1000$.

\begin{codeexample}
import numpy as np

def ddpm_sample(epsilon_theta, shape, betas, rng):
    """Sample one image from a trained DDPM noise predictor.
    epsilon_theta(x_t, t) -> predicted noise (same shape as x_t).
    betas: length-T variance schedule.
    """
    T          = len(betas)
    alphas     = 1.0 - betas
    alpha_bars = np.cumprod(alphas)
    x = rng.standard_normal(size=shape)        # x_T ~ N(0, I)
    for t in reversed(range(T)):
        eps_pred = epsilon_theta(x, t)
        coef1    = 1.0 / np.sqrt(alphas[t])
        coef2    = (1.0 - alphas[t]) / np.sqrt(1.0 - alpha_bars[t])
        mean     = coef1 * (x - coef2 * eps_pred)
        if t > 0:
            sigma = np.sqrt(betas[t])           # simple variance choice
            z     = rng.standard_normal(size=shape)
            x     = mean + sigma * z
        else:
            x     = mean                        # last step: no noise
    return x                                    # x_0 sample
\end{codeexample}

The loop is the algorithm. The schedule $\bar\alpha_t$ is precomputed once; the only per-step learned quantity is the noise prediction $\epsilon_\theta(x_t, t)$; the rest of the step is scalar arithmetic from the closed-form posterior above. Real implementations differ in three places: the schedule (linear, cosine, or learned), the variance choice $\sigma_t^2$, and the number of effective steps (the DDIM trick skips most of them by reusing the closed-form $q(x_t \mid x_0)$ identity rather than walking the original Markov chain). The core update is the same.

DDPM sampling costs $T$ forward passes through the noise predictor per generated sample, where each forward pass is an $O(\text{model FLOPs})$ U-Net or transformer evaluation. For $T = 1000$ and a model that runs in $1$ ms per pass on a modern GPU, that's about a second per image---two to three orders of magnitude slower than a single-pass VAE decoder or GAN generator. The accelerated samplers (DDIM, DPM-Solver, consistency models) trade quality for compute by reducing the effective $T$ to $10$--$50$ at the price of a small FID degradation. Training cost is dominated by the same forward-and-backward pass over many random $t$'s and is comparable to other large generative models; the asymmetry is concentrated at inference.

\subsection{Why Diffusion Models Work Well}

Diffusion models have several advantages over VAEs and GANs:

\begin{itemize}[leftmargin=*]
\item \emph{Training stability.} The training objective is a regression loss---predict the noise---with no adversarial dynamics. Optimization is often easier to monitor than GAN training, though architecture, noise schedule, data scale, and compute still matter.
\item \emph{Sample quality.} Decomposing generation into many small denoising steps lets diffusion models produce extremely detailed, high-quality outputs. Each step makes a small correction, sidestepping the single-shot generation challenge of GANs and VAEs.
\item \emph{Mode coverage.} Diffusion models are typically less prone to the classic GAN failure mode in which the generator collapses to a few sample types. Coverage is still an empirical question, though: rare cases, underrepresented slices, and conditioning failures must be checked directly.
\end{itemize}

The main disadvantage is computational cost. Generation requires running the denoising network for hundreds or thousands of steps---one forward pass per step---making it much slower than a single-pass GAN or VAE decoder.

Accelerated sampling methods such as DDIM and DPM-Solver reduce the number of denoising steps, trading some quality for speed \citep{song2020ddim,lu2022dpmsolver}. Tens of sampling steps often suffice in practice, compared with the original 1000-step schedule. During training, implementations typically sample timesteps from that schedule rather than unrolling every step for every example. For applications where generation speed matters---real-time augmentation during training, for example---this tradeoff is important.

\subsection{Guidance: Controlling What Gets Generated}

Diffusion models can be \emph{guided} toward specific outputs through a conditioning signal. \emph{Classifier-free guidance} trains the model both with and without the condition, then amplifies the difference at generation time \citep{ho2022classifierfree}. A guidance scale $w > 1$ pushes the output more strongly toward the condition:
\[
\hat{\epsilon} = \epsilon_\theta(x_t, \varnothing) + w \cdot \bigl(\epsilon_\theta(x_t, c) - \epsilon_\theta(x_t, \varnothing)\bigr)
\]

Higher guidance matches the condition more faithfully but produces less diversity. This is the \emph{quality--diversity tradeoff}: strong guidance gives you exactly what you asked for at the cost of repetitive outputs; weak guidance gives diverse outputs that may not match the condition well.

\begin{example}[Diffusion for Synthetic Training Data]
A team building a pedestrian detector needs training images of ``pedestrian crossing a wet road at dusk.'' A text-conditioned diffusion model generates 500 such images. With low guidance ($w=2$), the images are diverse---some show the road clearly, others are more abstract---but not all are realistic enough for training. With high guidance ($w=10$), the images are more uniform and realistic but lack diversity, with similar poses and similar lighting. The team settles on medium guidance ($w=5$) and post-filters the outputs, keeping 300 images that pass a visual quality check. These are then evaluated against real test frames to measure the sim-to-real gap.
\end{example}

\subsection*{Extension: Normalizing Flows and Flow Matching}

Two other generative approaches deserve mention; the rest of the book does not depend on this subsection. \emph{Normalizing flows} define an explicit, invertible mapping between a simple distribution (typically Gaussian) and the data distribution \citep{dinh2017realnvp}. Because the mapping is invertible, exact likelihood computation is possible---unlike VAEs, which use a lower bound, or GANs, which have no likelihood. The catch is that every transformation must be invertible with a tractable Jacobian, which limits architectural flexibility.

\emph{Flow matching} relaxes this constraint by defining a continuous-time flow from noise to data, trained with a regression objective related to diffusion and continuous-normalizing-flow views \citep{lipman2023flowmatching}. Sampling still requires numerically integrating the learned flow over multiple solver steps, but in practice fewer or cheaper steps than standard diffusion sampling can suffice while maintaining high sample quality. Flow matching is an active research area; treat it as a fast-moving family rather than a settled recipe.

Both methods share this chapter's theme: learning a mapping between a simple latent distribution and a complex data distribution, with the latent space providing structure for generation and interpolation.

\section{Evaluating Generative Models}
\label{sec:gen-eval}

Evaluating generative models is fundamentally harder than evaluating classifiers. A classifier's output can be compared to a ground-truth label. A generative model's output is a new data point with no single ``correct'' answer. Evaluation must assess two properties at once: \emph{quality} (do individual samples look realistic?) and \emph{diversity} (does the model cover the full range of variation in the data?).

\subsection{Automated Metrics}

\textbf{Fr\'{e}chet Inception Distance (FID).} A widely used metric for image generation \citep{heusel2017fid}. FID measures the distance between the distribution of real images and the distribution of generated images in the feature space of a pre-trained Inception network. Lower FID means the generated distribution is closer to the real one in that feature space. FID captures both quality (unrealistic images have unusual features) and diversity (a mode-collapsed generator has a narrower feature distribution), but only through the chosen representation.

\textbf{Inception Score (IS).} Measures two properties: each generated image should be confidently classified by an Inception network (quality), and the marginal distribution of predicted labels across generated images should have high entropy (diversity) \citep{salimans2016improved}. Higher IS is better under that proxy. IS has known limitations. It does not compare to the real data distribution, its label-diversity assumption may not match a single-domain dataset, and it can be gamed by producing one very realistic image per predicted class.

\textbf{Precision and recall for generation.} Adapted from classification: \emph{precision} estimates how much generated mass lies near the real data distribution (quality), and \emph{recall} estimates how much of the real data distribution is covered by generated samples (diversity) \citep{kynkaanniemi2019pr}. A mode-collapsed GAN can have high precision but low recall.

\begin{warningbox}
All automated metrics for generative models are proxies. FID depends on the choice of feature extractor---Inception was trained on ImageNet, not on your domain. IS does not compare to the real distribution. Precision-recall requires approximate support estimation in feature space, so its value depends on the representation and neighborhood method used. No single number captures whether generated data is ``good enough'' for a given downstream task. Always complement automated metrics with visual inspection and, when possible, downstream task performance.
\end{warningbox}

\subsection{Human Evaluation}

For many applications, human evaluation remains essential. Common protocols include:

\begin{itemize}[leftmargin=*]
\item \emph{Turing test / real-vs-fake.} Raters see a mix of real and generated samples and identify which are real. The fooling rate measures quality.
\item \emph{Preference ranking.} Raters compare outputs from two or more models and state which they prefer. Useful for comparing models, but it does not measure absolute quality.
\item \emph{Task-specific evaluation.} When generated data is used for augmentation, the downstream task metric---classification accuracy on real test data, for example---is the ultimate evaluation. A model that produces beautiful images that do not help the downstream task is not useful.
\end{itemize}

When generated data is used for training---data augmentation or synthetic training---the only evaluation that ultimately matters is downstream performance on a held-out \emph{real} test set. FID and IS are useful during development for comparing model variants, but the decision to ship should hinge on whether the generated data actually improves the system it supports.

A useful claim audit for generated pedestrian scenes runs as follows. \emph{Claim}: synthetic dusk scenes improve the real pedestrian detector under low-light conditions. \emph{Evidence available}: detector performance on a real held-out dusk slice before and after synthetic augmentation. \emph{Missing evidence}: whether gains persist across cameras, weather, and rare occlusions. \emph{Decision}: use synthetic data only if real-slice performance improves without degrading daytime or high-risk cases.

\section{Chapter Artifact: The Generative Model Audit Card}
\label{sec:gen-audit}

The chapter artifact is the \emph{generative model audit card}: a structured document recording what the model generates, how it was evaluated, and what failure modes to watch for. It complements the uncertainty audit card from Chapter~\ref{ch:probabilistic-models} and the data design card from Chapter~\ref{ch:dataset-design}.

\subsection{Generative Model Audit Card Structure}

\begin{table}[t]
\centering
\small
\setlength{\tabcolsep}{4pt}
\begin{tabular}{@{}L{2.8cm}L{7.2cm}@{}}
\toprule
Card field & What must be documented \\
\midrule
Model family & VAE, GAN, diffusion, flow, or hybrid; architecture details \\
Training data & What data the model was trained on; known gaps, biases, or limitations in the training set \\
Generation task & What the model is meant to generate and for what downstream purpose \\
Conditioning & What conditions or prompts control generation (class labels, text, images, none) \\
Quality metrics & FID, IS, precision/recall, or other automated metrics; values and what feature extractor was used \\
Diversity assessment & Evidence that the model covers the full range of variation; tests for mode collapse \\
Human evaluation & Protocol used, sample size, inter-rater agreement, and results \\
Known failure modes & What the model generates poorly (e.g., rare classes, unusual poses, specific lighting conditions) \\
Downstream impact & How generated data is used (augmentation, testing, creative); measured effect on the downstream task \\
Misuse risks & What harmful content the model could generate; what safeguards are in place \\
\bottomrule
\end{tabular}
\caption{A generative model audit card documents what the model generates, how it was evaluated, and what failure modes exist.}
\label{tab:gen-audit-card}
\end{table}

\subsection{Filled Example: Pedestrian Scene Generator}

\begin{example}[Generative Model Audit Card for Nighttime Pedestrian Augmentation]

\textbf{Model family.} Conditional diffusion model (U-Net backbone, 256$\times$256 resolution). Classifier-free guidance with scale $w = 5$.

\textbf{Training data.} 50,000 frames from campus shuttle cameras (72\% daytime, 25\% dusk/dawn, 3\% nighttime). Known gap: nighttime frames are underrepresented, which is the motivation for generation.

\textbf{Generation task.} Generate synthetic nighttime pedestrian scenes to augment the training set for a pedestrian detector.

\textbf{Conditioning.} Time-of-day label (day, dusk, night) and pedestrian count (0, 1, 2+).

\textbf{Quality metrics.} FID = 28.3 (all conditions), FID = 34.7 (nighttime only, compared to real nighttime frames). Feature extractor: Inception-v3 fine-tuned on campus imagery.

\textbf{Diversity assessment.} Visual inspection of 200 generated nighttime samples shows variation in pedestrian position, pose, and street lighting. All generated scenes, however, use the same two camera viewpoints from the training data. Rare events---wheelchair users, cyclists---are absent from generated nighttime scenes.

\textbf{Human evaluation.} Three annotators reviewed 100 generated nighttime scenes alongside 100 real ones. Fooling rate: 68\% (annotators labeled 68\% of generated scenes as ``real''). Inter-rater agreement (Fleiss's kappa): 0.54.

\textbf{Known failure modes.} Generated pedestrians sometimes have distorted limbs at close range. Rain effects are unconvincing because the training data contains very few rainy nighttime frames. Text on signs and license plates is garbled.

\textbf{Downstream impact.} Adding 2,000 generated nighttime scenes to the pedestrian detector's training set improved nighttime recall from 0.62 to 0.74 on the real nighttime test set. Daytime performance was unchanged ($\pm$0.01).

\textbf{Misuse risks.} The model can generate realistic campus scenes with identifiable buildings. Safeguard: model weights are stored on a restricted server, generation is logged, and generated images are watermarked for internal use only.

\end{example}

\section{Comparing the Three Families}

Each generative model family has distinct strengths. The right choice depends on the task, the data, and the computational budget.

\begin{table}[t]
\centering
\small
\setlength{\tabcolsep}{2pt}
\begin{tabular}{@{}L{2.05cm}L{2.58cm}L{2.58cm}L{2.58cm}@{}}
\toprule
& VAE & GAN & Diffusion \\
\midrule
Sample quality & Moderate; often blurry for images & High when training succeeds & Often very high in image and media domains when data and compute are sufficient \\
Diversity & Usually broad but decoder-dependent & Risk of mode collapse & Often strong, still needs coverage checks \\
Training stability & Usually easier to monitor & Unstable; needs tuning & Usually easier to monitor than GANs, but still sensitive to architecture, sampler, and compute \\
Generation speed & Fast (single pass) & Fast (single pass) & Slower; many solver steps unless accelerated \\
Likelihood & Lower bound (ELBO) & No likelihood & Usually approximate or estimator-dependent \\
Best for & Latent interpolation, anomaly detection, augmentation & Sharp image generation, style transfer & High-fidelity image or media generation when compute permits \\
\bottomrule
\end{tabular}
\caption{Trade-offs across the three main generative model families.}
\label{tab:gen-model-comparison}
\end{table}

\begin{decisionbox}
\textbf{Choosing a generative model:}
\begin{itemize}
\item If you need a structured latent space for interpolation, anomaly detection, or understanding data structure, start with a VAE.
\item If you need sharp, realistic samples and can invest in training stability, consider a GAN.
\item If sample quality is paramount and generation speed is not a bottleneck, consider a diffusion model and measure whether its diversity and conditioning behavior match the task.
\item If you need exact likelihoods (e.g., for density estimation or anomaly scoring), consider normalizing flows.
\item In all cases, evaluate on the downstream task, not just on sample quality metrics.
\end{itemize}
\end{decisionbox}

\section{Common Mistakes in Generative Modeling}

\textbf{Evaluating quality and ignoring diversity.} A model that produces a hundred beautiful but nearly identical samples is useless for data augmentation and dangerous as a basis for downstream training. Sample quality without a diversity audit hides mode collapse behind a flattering grid.

\textbf{Trusting FID alone.} FID depends on the feature extractor, the sample size, and the data domain, and small differences can flip rankings between extractors. A low FID is a proxy, not a guarantee that generated data will help on the downstream task.

\textbf{Letting generated data contaminate evaluation.} If synthetic samples augment the training set, the test set must stay real and untouched. Mixing generated and real data into evaluation turns the model into the judge of its own output.

\textbf{Ignoring the model's failure modes.} If the generator cannot produce rare events or underrepresented conditions, using its output as training data will not fix the coverage gap it was meant to address---and may amplify it. Deploying generated content without considering misuse compounds the problem; the generative model audit card exists to force these questions before release.

\section{Looking Ahead}

This chapter introduced three families of generative models and framed generation as a design decision. The choice of model, the conditioning scheme, and the evaluation protocol together shape what generated data can and cannot contribute to the downstream system. The generative model audit card documents these choices.

This chapter also closes Part~II. Across neural networks, probabilistic models, representation reuse, and generative representations, the unifying problem has been the same: learn structure from a static dataset and report honestly what the resulting model knows and does not know. The next part marks a category shift. In Part~III, ``Learning from Interaction,'' Chapter~\ref{ch:reinforcementlearning} stops asking what data we should collect and starts asking what the agent should \emph{do}---how a system can learn by acting, observing consequences, and updating its policy. Reward design replaces label design, exploration joins evaluation as a design constraint, and the data distribution is shaped by the agent itself. The book then returns to predictive learning across modalities in Part~IV with the interaction lens still active.

\section*{Chapter Summary}

Generative models learn to produce new data from latent representations. Variational autoencoders impose smooth probabilistic structure on the latent space, enabling interpolation and anomaly detection but producing blurrier samples. Generative adversarial networks use adversarial training to produce sharp samples but suffer from instability and mode collapse. Diffusion models decompose generation into iterative denoising, achieving the highest sample quality at the cost of slow generation. Evaluating generative models requires assessing both quality and diversity; automated metrics like FID and IS are useful proxies but not definitive. Generated data should always be validated against real held-out data. The generative model audit card documents what the model generates, how it was evaluated, what failure modes exist, and what misuse risks are present.

\begin{studyquestions}
\item How does a VAE's latent space differ from a standard autoencoder's latent space?
\item What do the two terms in the VAE loss (reconstruction and KL divergence) each encourage?
\item Why does the reparameterization trick enable backpropagation through sampling?
\item Why do GANs tend to produce sharper images than VAEs?
\item What is mode collapse, and why is it hard to detect with standard training metrics?
\item How do diffusion models decompose generation into many small steps, and why does this help quality?
\item What is the quality--diversity tradeoff in classifier-free guidance?
\item Why is FID an imperfect measure of generative quality?
\item When is downstream task performance a better evaluation metric than FID or IS?
\item What should a generative model audit card document about failure modes?
\end{studyquestions}

\begin{exercises}
\item \exercisetag{Calculation}{Core}Compute the closed-form Gaussian KL by hand. For $q(z) = \mathcal{N}(\mu, \sigma^2)$ and $p(z) = \mathcal{N}(0, 1)$ in one dimension, the formula $D_\text{KL}(q \,\|\, p) = \tfrac{1}{2}(\mu^2 + \sigma^2 - \log \sigma^2 - 1)$ should give exactly $0$ when $\mu = 0$ and $\sigma^2 = 1$. Verify this by substitution. Then compute the KL for $(\mu, \sigma^2) = (0, 4)$ and $(\mu, \sigma^2) = (3, 1)$, and explain in one sentence which of the two posteriors is being penalized more by the KL regularizer and why.
\item \exercisetag{Proof}{Advanced}Re-derive the ELBO from the marginal log-likelihood. Starting from $\log p(x) = \log \int p(x \mid z) p(z)\,dz$, multiply and divide the integrand by $q(z \mid x)$, apply Jensen's inequality (recall $\log$ is concave), and split the resulting lower bound into the reconstruction and KL-regularization terms. Then identify the slack between $\log p(x)$ and the ELBO as a KL divergence between $q(z \mid x)$ and the true posterior $p(z \mid x)$.
\item \exercisetag{Implementation}{Advanced}Train a VAE on MNIST or Fashion-MNIST. Visualize the latent space (using a 2D latent code) by coloring points by their class label. Do classes form distinct regions? Sample new digits by decoding random latent points and inspect the results.
\item \exercisetag{Implementation}{Intermediate}Using the trained VAE from the previous exercise, interpolate between two digits in latent space (e.g., from a ``3'' to a ``7''). Show the decoded images at 5 equally spaced points along the interpolation path. Is the transition smooth?
\item \exercisetag{Diagnosis}{Intermediate}Given two provided sample grids, one from a VAE-like model and one from a GAN-like model, assess: (a) visual quality of individual samples, (b) diversity across the samples. Which model appears sharper? Which appears more diverse? State what evidence the sample grids cannot provide.
\item \exercisetag{Diagnosis}{Intermediate}Given two reported FID scores computed with different feature extractors, decide whether they support the same ranking of two generators. What does disagreement between the scores tell you about FID's reliability?
\item \exercisetag{Design}{Intermediate}Design a data augmentation experiment using a generative model of your choice. Specify: the downstream task, the generative model, the conditioning scheme, the number of generated samples, and the evaluation protocol. Explicitly state what the real held-out test set is and why the generated data must not contaminate it.
\item \exercisetag{Design}{Advanced}Construct a generative model audit card (Table~\ref{tab:gen-audit-card}) for a generative model you have trained or studied. Include at least two known failure modes and one misuse risk.
\end{exercises}

\begin{challengeexercises}
\item \exercisetag{Diagnosis}{Advanced}A team uses a GAN to augment pedestrian detection training data. After augmentation, overall detection accuracy improves by 5\%, but detection of wheelchair users gets worse. Diagnose the problem and propose a fix.
\item \exercisetag{Diagnosis}{Advanced}Explain why a diffusion model with 1000 denoising steps produces higher-quality samples than one with 20 steps, but why 20 steps may be sufficient for a given application. Under what conditions would you choose fewer steps?
\item \exercisetag{Diagnosis}{Advanced}A VAE trained on student messages has a well-organized latent space where nearby points decode to similar messages. But when you sample from the prior $\mathcal{N}(0, I)$, many decoded messages are incoherent. Explain what has gone wrong (hint: consider the gap between the aggregate posterior and the prior) and propose a mitigation.
\item \exercisetag{Design}{Advanced}Design an evaluation protocol for a text-conditioned diffusion model used to generate synthetic medical images. What metrics would you use? What human evaluation would you require? What regulatory or ethical considerations apply?
\end{challengeexercises}
